% Lecture 10, page 2 -- continuation of curs10-pagina1-en.tex.
% Fragment for \input; no preamble, \clearpage or new lecture heading.
\subsection*{Lecture 10 - Page 2}

\bigskip
\noindent\underline{Ex.} \textbf{3)} Let $G$ be an arbitrary group,
$K,H\leq G$, \quad $A=\{aH\mid a\in G\}$ -- the set of left cosets
of $H$ in $G$.

\bigskip
\noindent $K\times A\xrightarrow{\ \psi\ }A$.

\noindent $(k,aH)\overset{\psi}{\longmapsto}k\cdot aH=(ka)H$,
\quad $k\in K$, $aH\in A$.

\bigskip
\noindent $\psi$ -- the action by left translation on the set of cosets
of $H$.

\bigskip
\noindent If $A$ is a $G$-set, we may associate with it the
 equivalence relation ``$\sim_G$''.

\medskip
\noindent $x,y\in A$, \quad
$x\sim_G y\ \Longleftrightarrow\ (\exists)\,a\in G$
such that $ax=y$.

\bigskip
\noindent If $x\in A$, \quad
$\widehat{x}\overset{\mathrm{def}}{=}\operatorname{orb}_G(x)
=Gx=\{ax\mid a\in G\}$ -- \underline{the orbit of $x$}.

\bigskip
\noindent\underline{The stabiliser of $x$ in $G$}:

\medskip
\noindent $\operatorname{Stab}_G(x)=\{a\in G\mid ax=x\}\leq G$.

\bigskip
\noindent If $x\in A$, \quad
$\bigl|\operatorname{orb}_G(x)\bigr|
=[G:\operatorname{Stab}_G(x)]$.

\medskip
\noindent If $|G|<\infty$, \quad
$\bigl|\operatorname{orb}_G(x)\bigr|\mid |G|$
(by Lagrange's theorem).

\bigskip
\noindent\underline{Ex. (1)}:
\quad $\operatorname{Stab}_G(x)\overset{\mathrm{def}}{=}
\{a\in G\mid axa^{-1}=x\}
=\{a\in G\mid ax=xa\}$.

\medskip
\noindent In this case,
$\operatorname{Stab}_G(x)\overset{\mathrm{not.}}{=}C_G(x)$
-- \underline{the centraliser of $x$ in $G$}.

\medskip
\noindent $\Longrightarrow\ \bigl|\operatorname{orb}_G(x)\bigr|
=[G:C_G(x)]$;
\quad $G$ commutative $\Longrightarrow C_G(x)=G$.

\bigskip
\noindent\underline{Ex. (2)}: \quad $H\leq G$.

\medskip
\noindent $\operatorname{Stab}_G(H)\overset{\mathrm{def}}{=}
\{a\in G\mid a\cdot H=H\}
=\{a\in G\mid aHa^{-1}=H\}$.

\medskip
\noindent We write $N_G(H)$ -- \underline{the normaliser of $H$ in $G$}.

\bigskip
\noindent\underline{Prop.}
\begin{enumerate}
\item[(1)] $H\subseteq N_G(H)$.
\item[(2)] $H\trianglelefteq N_G(H)$.
\item[(3)] $N_G(H)$ is the largest subgroup of $G$
 in which $H$ is a normal subgroup.
\end{enumerate}

\medskip
\noindent\underline{Proof.}

\noindent (1) \quad $h\in H$.

\noindent (and $a\in H$).\footnote{The last parenthetical phrase in the manuscript is difficult to read; ``and $a\in H$'' is provisional. The proof continues on the following page.}
